\section{Deterministic certificate and finite stress tests}
\label{sec:certificate}

The release package separates theorem proof from executable reproduction.
The producer records the source object, exact rational exponent ledger,
theorem status, route firewall, and a finite algebraic row fixture in canonical
JSON.  The leading rational $1197/800$, the powers $1/96$ and
$\pm23/2400$, and the interval $[1/2,1]$ are represented by exact rational
pairs before serialization.  A SHA-256 digest binds the canonical payload.

The independent checker does not import the producer.  It reconstructs the
entire certificate and rejects nine mutation classes: status promotion,
Boolean/integer confusion, the wrong rational constant, the wrong exponent,
a coefficient-sign change, a profile-class change, expansion of the
top-prime domain, substitution of a q-collapsed object, and substitution of an
inadmissible plateau.

A separate stress program uses two functions of the form
\[
 \psi(t)=\frac{\exp(-1/(1-t^2))w(t)}
 {\int_{-1}^{1}\exp(-1/(1-s^2))w(s)\,ds}
 \one_{|t|<1},
\]
with $w(t)=1$ and $w(t)=1+t^2/4$.  These are fixed, nonnegative smooth
profiles satisfying the literal support condition; numerical quadrature checks
their normalization and bound.  Across three increasing prime fixtures, the
program checks residue injectivity, equality of atom and residue-row energies,
and decreasing error in the approximation of $\kappa_\psi$.

All finite values in this section are classified
\texttt{NUMERICAL\_FINITE\_ILLUSTRATION\_ONLY}.  They test implementation
and schema integrity but do not support the asymptotic theorem.  The proof of
Theorem~\ref{thm:main} rests on Lemmas~\ref{lem:row} and~\ref{lem:riemann}
and the classical weighted prime number theorem.
